%% kapa: exam=1 round=31 year=2020 period=1 form=A subject=civil_law subject_ko=민법 seq=39 number=39
%% source: Q-Net 감정평가사 기출문제 (cst00309) artlSeq=5207219 file=extracted/1008/5207219/A형/A형/2020년도 제31회 감정평가사 1차 시험_1교시_A형.pdf
%% conversion: kordoc 4.24.51 -> tools/gen1.py (auto); review=proofread-sonnet; flags=none
\begin{question}{39}
저당권에 관한 설명으로 옳은 것은? (다툼이 있으면 판례에 따름)
\begin{choices}
  \item 저당부동산의 소유권이 제3자에게 양도된 후 피담보채권이 변제된 때에는 저당권을 설정한 종전소유자도 저당권설정등기의 말소를 청구할 권리가 있다.
  \item 저당권을 설정한 사람이 채무자가 아닌 경우, 그는 원본채권이 이행기를 경과한 때부터 1년분의 범위에서 지연배상을 변제할 책임이 있다.
  \item 근저당권의 채무자가 피담보채권의 일부를 변제한 경우, 변제한 만큼 채권최고액이 축소된다.
  \item 저당권자는 배당기일 전까지 물상대위권을 행사하여 우선변제를 받을 수 있다.
  \item 대체물채권을 담보하기 위하여 저당권을 설정한 경우, 피담보채권액은 채권을 이행할 때의 시가로 산정한 금액으로 한다.
\end{choices}
\end{question}
\answer{1}
